% Pillar-1 (scaling) standalone abstract
\begin{abstract}
Group Relative Policy Optimization (GRPO) has become a default algorithm for
reinforcement-learning (RL) post-training of language models, yet how its benefit
scales with model size remains unclear and is easily confounded by implementation
and evaluation choices. We study GRPO scaling as one pillar of \ours{}, a
benchmark spanning 70+ runs across seven RL libraries and five model families
($0.6$B--$\sim$685B parameters) on GSM8K, HumanEval, and synthetic tool-use tasks.
Our central result is a conservative negative one: across roughly $2.4$ orders of magnitude
in parameter count across the extended anchor pool, the cross-scale slope of the mean GRPO training reward is statistically
indistinguishable from zero, and no single saturation exponent is identifiable
(four of the five frontier-scale traces have their saturation rate pinned at the
fit boundary, and none is identifiable). A pre-registered test geometrically falsifies a clean three-phase
saturation hypothesis (only $2$ of $12$ anchors match), and the Nemotron-120B run
is better described as a distinct collapse phase (zero-reward step fraction $0.55$
versus $\le 0.067$ elsewhere) than as a point on a smooth trend. The defensible
object is therefore a local, stack-conditioned taxonomy of endpoint ceilings,
with categorical capability structure carrying more signal than $\log_{10} N$,
not a Chinchilla-style power law in $N$. We
release all code, logs, and checkpoints.
\end{abstract}
